% GENERATED FILE. Edit canonical lessons, then run npm run generate:books.
% canonical-source-hash: fnv1a64:382d54e971142744
% canonical-lessons: JA-C117-hotate, JA-C117-nomi, JA-C117-koumori, JA-C117-kamoshika, JA-C117-suitou

\chapter{Scallop, Flea, Bat, Japanese serow, Water flask}
\label{ch:ja-scallop-flea-bat-japanese-serow-water-flask}
\hlchaptermodality{117}

\begin{chapteropening}
\textbf{By the end of this chapter:} \emph{I can say a scallop, a flea, a bat, a Japanese serow and a water flask.}

Complete the last lesson of chapter 117: I can say a scallop, a flea, a bat, a Japanese serow and a water flask.
\end{chapteropening}

\section[\texorpdfstring{\ja{ほたて} (hotate)}{ほたて (hotate)}]{\ja{ほたて} (hotate) — a scallop}

\label{lesson:JA-C117-hotate}



\begin{quote}
Before the new one: say the Japanese for a Pacific saury, then the Japanese for a short-neck clam.
\end{quote}

\subsection*{You'll want to know: \ja{ほたて}}
\textbf{\ja{ほたて}} — \emph{hotate} — \textquotedblleft{}a scallop\textquotedblright{}.

\begin{grammarlens}[title={What you've built}]
One more word for this chapter.
\end{grammarlens}

\subsection*{Your turn}
\begin{itemize}
\raggedright
  \item \emph{Say it:} \emph{hotate}
  \item \emph{Say it:} \emph{hotate}, once more
  \item \emph{Say it:} say \emph{hotate}
  \item \emph{Recall:} say the Japanese for a Pacific saury, then the Japanese for a short-neck clam, then say \emph{hotate} again
\end{itemize}

\begin{tcolorbox}[breakable,colback=teal!4,colframe=teal!35!black,title={Before you move on}]
What does \ja{ほたて} mean? (\textquotedblleft{}A scallop\textquotedblright{}.) Say it once more.
\end{tcolorbox}
\section[\texorpdfstring{\ja{のみ} (nomi)}{のみ (nomi)}]{\ja{のみ} (nomi) — a flea}

\label{lesson:JA-C117-nomi}



\begin{quote}
Before the new one: say the Japanese for a short-neck clam, then the Japanese for a scallop.
\end{quote}

\subsection*{You'll want to know: \ja{のみ}}
\textbf{\ja{のみ}} — \emph{nomi} — \textquotedblleft{}a flea\textquotedblright{}.

\begin{grammarlens}[title={What you've built}]
One more word for this chapter.
\end{grammarlens}

\subsection*{Your turn}
\begin{itemize}
\raggedright
  \item \emph{Say it:} \emph{nomi}
  \item \emph{Say it:} \emph{nomi}, once more
  \item \emph{Say it:} say \emph{nomi}
  \item \emph{Recall:} say the Japanese for a short-neck clam, then the Japanese for a scallop, then say \emph{nomi} again
\end{itemize}

\begin{tcolorbox}[breakable,colback=teal!4,colframe=teal!35!black,title={Before you move on}]
What does \ja{のみ} mean? (\textquotedblleft{}A flea\textquotedblright{}.) Say it once more.
\end{tcolorbox}
\section[\texorpdfstring{\ja{こうもり} (koumori)}{こうもり (koumori)}]{\ja{こうもり} (koumori) — a bat}

\label{lesson:JA-C117-koumori}



\begin{quote}
Before the new one: say the Japanese for a scallop, then the Japanese for a flea.
\end{quote}

\subsection*{You'll want to know: \ja{こうもり}}
\textbf{\ja{こうもり}} — \emph{koumori} — \textquotedblleft{}a bat\textquotedblright{}.

\begin{grammarlens}[title={What you've built}]
One more word for this chapter.
\end{grammarlens}

\subsection*{Your turn}
\begin{itemize}
\raggedright
  \item \emph{Say it:} \emph{koumori}
  \item \emph{Say it:} \emph{koumori}, once more
  \item \emph{Say it:} say \emph{koumori}
  \item \emph{Recall:} say the Japanese for a scallop, then the Japanese for a flea, then say \emph{koumori} again
\end{itemize}

\begin{tcolorbox}[breakable,colback=teal!4,colframe=teal!35!black,title={Before you move on}]
What does \ja{こうもり} mean? (\textquotedblleft{}A bat\textquotedblright{}.) Say it once more.
\end{tcolorbox}
\section[\texorpdfstring{\ja{かもしか} (kamoshika)}{かもしか (kamoshika)}]{\ja{かもしか} (kamoshika) — a Japanese serow}

\label{lesson:JA-C117-kamoshika}



\begin{quote}
Before the new one: say the Japanese for a flea, then the Japanese for a bat.
\end{quote}

\subsection*{You'll want to know: \ja{かもしか}}
\textbf{\ja{かもしか}} — \emph{kamoshika} — \textquotedblleft{}a Japanese serow\textquotedblright{}.

\begin{grammarlens}[title={What you've built}]
One more word for this chapter.
\end{grammarlens}

\subsection*{Your turn}
\begin{itemize}
\raggedright
  \item \emph{Say it:} \emph{kamoshika}
  \item \emph{Say it:} \emph{kamoshika}, once more
  \item \emph{Say it:} say \emph{kamoshika}
  \item \emph{Recall:} say the Japanese for a flea, then the Japanese for a bat, then say \emph{kamoshika} again
\end{itemize}

\begin{tcolorbox}[breakable,colback=teal!4,colframe=teal!35!black,title={Before you move on}]
What does \ja{かもしか} mean? (\textquotedblleft{}A Japanese serow\textquotedblright{}.) Say it once more.
\end{tcolorbox}
\section[\texorpdfstring{\ja{すいとう} (suitou)}{すいとう (suitou)}]{\ja{すいとう} (suitou) — a water flask}

\label{lesson:JA-C117-suitou}



\begin{quote}
Before the new one: say the Japanese for a bat, then the Japanese for a Japanese serow.
\end{quote}

\subsection*{You'll want to know: \ja{すいとう}}
\textbf{\ja{すいとう}} — \emph{suitou} — \textquotedblleft{}a water flask\textquotedblright{}.

\begin{grammarlens}[title={What you've built}]
One more word for this chapter.
\end{grammarlens}

\subsection*{Your turn}
\begin{itemize}
\raggedright
  \item \emph{Say it:} \emph{suitou}
  \item \emph{Say it:} \emph{suitou}, once more
  \item \emph{Say it:} say \emph{suitou}
  \item \emph{Recall:} say the Japanese for a bat, then the Japanese for a Japanese serow, then say \emph{suitou} again
\end{itemize}

\begin{tcolorbox}[breakable,colback=teal!4,colframe=teal!35!black,title={Before you move on}]
What does \ja{すいとう} mean? (\textquotedblleft{}A water flask\textquotedblright{}.) Say it once more.
\end{tcolorbox}
